% from ega4-16-fr.tex
\oldpage[IV-4]{28}
d'homomorphismes d'Algèbres, de l'interprétation analogue donnée dans
(0, 20.1.6), et de la correspondance canonique entre homomorphismes de
$\mathcal{O}_X$-Algèbres et homomorphismes de $B$-algèbres (I, 1.3.13 et
1.3.8).

\begin{proposition}[16.5.3]
\label{IV.16.5.3}
Soit $f=(\psi,\theta):X\to S$ un morphisme de préschémas.
\begin{enumerate}
\item[(i)] La différentielle
$d_{X/S}:\mathcal{O}_X\to\Omega_{X/S}^1$ (16.3.6) est une $S$-dérivation.
\item[(ii)] Pour tout $\mathcal{O}_X$-Module $\mathscr{F}$, l'application
$u\mapsto u\circ d_{X/S}$ est un isomorphisme de
$\Gamma(X,\mathcal{O}_X)$-modules
\[
  \operatorname{Hom}_{\mathcal{O}_X}(\Omega_{X/S}^1,\mathscr{F})
  \xrightarrow{\sim}
  \operatorname{D\acute{e}r}_S(\mathcal{O}_X,\mathscr{F}).
  \tag{16.5.3.1}
  \label{IV.16.5.3.1}
\]
\end{enumerate}
\end{proposition}

L'assertion (i) a déjà été notée (16.3.6). D'autre part, il est immédiat (en
vertu de (0, 20.4.8)) que $u\mapsto u\circ d_{X/S}$ est injectif, en
considérant les restrictions à une fibre $\mathcal{O}_x$ des deux membres et
utilisant (16.4.16.2). Pour voir que l'homomorphisme (16.5.3.1) est surjectif,
considérons une $S$-dérivation $D:\mathcal{O}_X\to\mathscr{F}$; pour tout
ouvert affine $V=\operatorname{Spec}(B)$ de $X$, tel que $f(V)$ soit contenu
dans un ouvert affine $U=\operatorname{Spec}(A)$ de $S$,
$D_V:B\to\Gamma(V,\mathscr{F})$ est une $A$-dérivation, donc il existe un
\emph{unique} $B$-homomorphisme
$u_V:\Omega_{B/A}^1\to\Gamma(V,\mathscr{F})$ tel que
$D_V=u_V\circ d_{B/A}$ (0, 20.4.8); en outre l'unicité de $u_V$ montre
aussitôt que pour tout ouvert affine $W\subset V$, on a $u_W=u_V|W$, donc les
$u_V$ définissent un homomorphisme de $\mathcal{O}_X$-Modules
$u:\mathcal{O}_X\to\mathscr{F}$ répondant à la question.

\begin{env}[16.5.4]
\EGAReaderTarget{ega4.c04.statement.16-5-4}{IV.16.5.4}
Avec les notations de (16.5.1), pour tout ouvert $U$ de $X$,
$\operatorname{D\acute{e}r}_S(\mathcal{O}_U,\mathscr{F}|U)$ est un
$\Gamma(U,\mathcal{O}_X)$-module et il est clair que l'application
$U\mapsto\operatorname{D\acute{e}r}_S(\mathcal{O}_U,\mathscr{F}|U)$ est un
préfaisceau; en fait, c'est même un \emph{faisceau} (donc un
$\mathcal{O}_X$-Module), en vertu de la caractérisation ponctuelle des
$S$-dérivations, vue en (16.5.1). Cet $\mathcal{O}_X$-Module se note
$\mathscr{D}\!\mathrm{\acute{e}r}_S(\mathcal{O}_X,\mathscr{F})$ et est appelé
\emph{faisceau des $S$-dérivations de $\mathcal{O}_X$ dans $\mathscr{F}$}, et
ce qu'on vient de voir s'exprime encore par le corollaire suivant:
\end{env}

\begin{corollary}[16.5.5]
\label{IV.16.5.5}
Pour tout $\mathcal{O}_X$-Module $\mathscr{F}$, l'homomorphisme de
$\mathcal{O}_X$-Modules déduit de $u\mapsto u\circ d_{X/S}$
\[
  \mathscr{H}\!om_{\mathcal{O}_X}(\Omega_{X/S}^1,\mathscr{F})
  \longrightarrow
  \mathscr{D}\!\mathrm{\acute{e}r}_S(\mathcal{O}_X,\mathscr{F})
  \tag{16.5.5.1}
  \label{IV.16.5.5.1}
\]
est bijectif.
\end{corollary}

\begin{corollary}[16.5.6]
\label{IV.16.5.6}
\begin{enumerate}
\item[(i)] Si le morphisme $f:X\to S$ est localement de présentation finie et
si $\mathscr{F}$ est un $\mathcal{O}_X$-Module quasi-cohérent, alors
$\mathscr{D}\!\mathrm{\acute{e}r}_S(\mathcal{O}_X,\mathscr{F})$ est un
$\mathcal{O}_X$-Module quasi-cohérent.
\item[(ii)] Si de plus $S$ est localement noethérien et si $\mathscr{F}$ est
cohérent, alors $\mathscr{D}\!\mathrm{\acute{e}r}_S(\mathcal{O}_X,\mathscr{F})$ est un
$\mathcal{O}_X$-Module cohérent.
\end{enumerate}
\end{corollary}

L'assertion (i) résulte de l'isomorphisme (16.5.5.1), de (16.4.22) et
(I, 1.3.12); l'assertion (ii) résulte de (0$_{\rm I}$, 5.3.5).

\begin{env}[16.5.7]
\EGAReaderTarget{ega4.c04.statement.16-5-7}{IV.16.5.7}
On pose
\[
  \mathfrak{G}_{X/S}=
  \mathscr{H}\!om_{\mathcal{O}_X}(\Omega_{X/S}^1,\mathcal{O}_X)
  =\mathscr{D}\!\mathrm{\acute{e}r}_S(\mathcal{O}_X,\mathcal{O}_X)
  \tag{16.5.7.1}
  \label{IV.16.5.7.1}
\]
et on dit que c'est le \emph{faisceau des $S$-dérivations de
$\mathcal{O}_X$}, ou encore le \emph{faisceau tangent de $X$ relativement à
$S$}: c'est donc le \emph{dual} du $\mathcal{O}_X$-Module
$\Omega_{X/S}^1$. Si $f$ est localement de présentation
